\documentclass[11pt,a4paper]{article}
\usepackage[utf8]{inputenc}
\usepackage[T1]{fontenc}
\usepackage{geometry}
\geometry{margin=1in}
\usepackage{hyperref}
\usepackage{enumitem}
\usepackage{booktabs}
\usepackage{longtable}
\usepackage{xcolor}
\usepackage{titlesec}
\usepackage{amsmath,amssymb}

\titleformat{\section}{\Large\bfseries\color{blue!70!black}}{}{0em}{}
\titleformat{\subsection}{\large\bfseries\color{blue!50!black}}{}{0em}{}
\titleformat{\subsubsection}{\normalsize\bfseries\color{blue!30!black}}{}{0em}{}

\title{\textbf{Replication Plan}\\[0.5em]
\large Clustering Huge Protein Sequence Sets in Linear Time\\[0.3em]
\normalsize OSTI ID: 1624105}
\author{Auto-generated Replication Blueprint}
\date{\today}

\begin{document}
\maketitle
\tableofcontents
\newpage

%---------------------------------------------------------------------
\section{Paper Overview}
%---------------------------------------------------------------------
This paper presents Linclust, a linear-time clustering algorithm for protein sequence databases, implemented within the MMseqs2 software suite. Unlike CD-HIT and UCLUST which scale quadratically or worse, Linclust achieves $O(N)$ runtime by using a $k$-mer-based pre-filtering step to identify candidate cluster members, followed by alignment-based verification. The paper benchmarks Linclust on UniProt (hundreds of millions of sequences), demonstrating orders-of-magnitude speedup over competitors while maintaining comparable clustering quality.

%---------------------------------------------------------------------
\section{Detailed Workflow}
%---------------------------------------------------------------------

\subsection{Phase 1: Software Installation}
\begin{enumerate}[label=\textbf{1.\arabic*}]
  \item \textbf{Install MMseqs2} (includes Linclust) from GitHub or conda.
  \item \textbf{Install competitor tools} for benchmarking: CD-HIT, UCLUST/USEARCH, DIAMOND, MASH.
  \item Verify installations with small test datasets.
\end{enumerate}

\subsection{Phase 2: Dataset Acquisition}
\begin{enumerate}[label=\textbf{2.\arabic*}]
  \item \textbf{Download UniProt/UniRef datasets.}
    \begin{itemize}
      \item UniRef100, UniRef90, UniRef50 FASTA files.
      \item Specific release version as used in the paper (if identifiable).
    \end{itemize}
  \item \textbf{Prepare subsets} for scaling experiments (e.g., 1M, 10M, 50M, 100M, full database).
\end{enumerate}

\subsection{Phase 3: Clustering Experiments}
\begin{enumerate}[label=\textbf{3.\arabic*}]
  \item \textbf{Run Linclust} at multiple sequence identity thresholds (e.g., 30\%, 50\%, 70\%, 90\%).
    \begin{itemize}
      \item \texttt{mmseqs linclust inputDB outputDB tmp --min-seq-id 0.3}
      \item Record wall-clock time, peak memory usage, number of clusters.
    \end{itemize}
  \item \textbf{Run competitor tools} on the same datasets with equivalent parameters.
  \item \textbf{Scaling experiments.}
    \begin{itemize}
      \item Run on increasing dataset sizes; plot runtime vs.\ $N$.
      \item Verify linear scaling of Linclust vs.\ super-linear scaling of competitors.
    \end{itemize}
\end{enumerate}

\subsection{Phase 4: Quality Assessment}
\begin{enumerate}[label=\textbf{4.\arabic*}]
  \item \textbf{Cluster quality metrics.}
    \begin{itemize}
      \item Sensitivity: fraction of true homologs placed in the same cluster.
      \item Specificity: fraction of cluster members that are true homologs.
      \item Comparison against SCOP/SCOPe or Pfam domain assignments as ground truth.
    \end{itemize}
  \item \textbf{Cluster size distribution} analysis.
  \item \textbf{Alignment quality}: sample and verify pairwise alignments within clusters.
\end{enumerate}

\subsection{Phase 5: Figure and Table Reproduction}
\begin{enumerate}[label=\textbf{5.\arabic*}]
  \item Reproduce runtime vs.\ database-size scaling plots.
  \item Reproduce quality comparison tables.
  \item Reproduce cluster size distribution histograms.
\end{enumerate}

%---------------------------------------------------------------------
\section{Datasets}
%---------------------------------------------------------------------

\begin{longtable}{p{4.5cm} p{3.5cm} p{6cm}}
\toprule
\textbf{Dataset / Source} & \textbf{Type} & \textbf{Location / Notes} \\
\midrule
\endfirsthead
\toprule
\textbf{Dataset / Source} & \textbf{Type} & \textbf{Location / Notes} \\
\midrule
\endhead
UniProt / UniRef & Public protein sequences & \url{https://www.uniprot.org/downloads} \\
\midrule
SCOP / SCOPe & Structural classification & \url{https://scop.berkeley.edu/} \\
\midrule
Pfam & Protein family database & \url{https://www.ebi.ac.uk/interpro/} (merged into InterPro) \\
\midrule
Metaclust & Metagenomic protein clusters & \url{https://metaclust.mmseqs.com/} \\
\bottomrule
\end{longtable}

%---------------------------------------------------------------------
\section{Models, Tools, and Software}
%---------------------------------------------------------------------

\subsection{Core Tools}
\begin{longtable}{p{4.5cm} p{2.5cm} p{6.5cm}}
\toprule
\textbf{Tool / Library} & \textbf{Version} & \textbf{Purpose / URL} \\
\midrule
\endfirsthead
\toprule
\textbf{Tool / Library} & \textbf{Version} & \textbf{Purpose / URL} \\
\midrule
\endhead
MMseqs2 (includes Linclust) & latest & Primary clustering tool \newline \url{https://github.com/soedinglab/MMseqs2} \\
\midrule
CD-HIT & $\geq 4.8$ & Competitor clustering tool \newline \url{https://github.com/weizhongli/cdhit} \\
\midrule
USEARCH / UCLUST & latest & Competitor clustering \newline \url{https://www.drive5.com/usearch/} \\
\midrule
DIAMOND & $\geq 2.0$ & Fast protein alignment \newline \url{https://github.com/bbuchfink/diamond} \\
\midrule
MASH & latest & Fast genome/sequence distance \newline \url{https://github.com/marbl/Mash} \\
\midrule
Python / Matplotlib & $\geq 3.8$ & Benchmarking scripts and visualization \\
\midrule
GNU time / /usr/bin/time & system & Wall-clock time and memory measurement \\
\bottomrule
\end{longtable}

\subsection{Hardware Requirements}
\begin{itemize}
  \item \textbf{For full UniProt}: $\geq$ 128\,GB RAM, $\geq$ 500\,GB SSD storage.
  \item \textbf{For subsets (10M sequences)}: $\geq$ 32\,GB RAM.
  \item Multi-core CPU (MMseqs2 uses OpenMP for parallelism).
\end{itemize}

\end{document}
